\chapter{Analisis dan Solusi}

Fenomena ketidaksinkronan pada arsitektur CQRS muncul sebagai konsekuensi dari pemisahan model tulis dan baca yang bergantung pada aliran event.
Pada kondisi beban yang fluktuatif, laju kedatangan event tidak selalu sejalan dengan kapasitas pemrosesan konsumer, sehingga backlog dapat meningkat dan menghasilkan keterlambatan pembaruan pada database baca.
Peningkatan ukuran batch memang dapat menurunkan frekuensi pemrosesan dan mengurangi overhead, tetapi juga berpotensi menambah penundaan pembaruan data.
Sebaliknya, ukuran batch yang terlalu kecil menghasilkan pemrosesan yang lebih responsif, namun cenderung meningkatkan konsumsi sumber daya dan menurunkan throughput.

Untuk itu, diperlukan sistem yang mampu menyeimbangkan kebutuhan akan data yang responsif namun juga efisien dalam penggunaan sumber daya.
Dalam perancangan sistem tersebut, diperlukan pemahaman yang baik mengenai hubungan antara variabel-variabel yang membentuk dinamika sistem.
Analisis ini mencakup identifikasi variabel state yang mewakili kondisi sistem pada suatu waktu, variabel kontrol yang dapat diatur, serta bagaimana perubahan pada variabel tersebut memengaruhi performa sinkronisasi, baik dalam hal backlog, latensi, maupun utilisasi sumber daya.
Pemahaman ini menjadi dasar perumusan model pengendalian yang memungkinkan penyesuaian parameter secara adaptif agar sistem tetap berada pada kondisi yang stabil dan efisien.

\section{Analisis Masalah}

Alur sinkronisasi pada arsitektur CQRS bergantung pada proses konsumsi event dari log, pengelompokan event ke dalam batch, dan materialisasi ulang menjadi read model.
Pada kondisi beban rendah, alur ini cenderung stabil karena kapasitas pemrosesan sejalan dengan laju kedatangan event.
Namun ketika terjadi peningkatan throughput penulisan atau penurunan kapasitas konsumsi, backlog mulai meningkat dan menimbulkan keterlambatan propagasi data ke read model.
Keterlambatan ini dapat secara langsung dirasakan pengguna dalam hal ketidakakuratan data karena data yang tidak aktual.

Masalah menjadi lebih kompleks karena pengaturan parameter utama seperti ukuran batch dan interval polling tidak selalu menaikkan performa.
Ukuran batch yang lebih besar dapat meningkatkan efisiensi pemrosesan karena mengurangi frekuensi commit dan overhead jaringan dan komputasi, tetapi juga menunda waktu mulai pemrosesan sehingga menambah latensi sinkronisasi terlihat pada Gambar~\ref{fig:batch-latency}.
Sebaliknya, ukuran batch yang lebih kecil mempercepat pembaruan data tetapi meningkatkan konsumsi sumber daya CPU dan I/O, yang pada tingkat tertentu justru menurunkan throughput keseluruhan.
Ketidakseimbangan ini menyebabkan sistem rentan terhadap ketidakpastian performa ketika beban berubah-ubah, terutama pada skenario bersifat \emph{bursty}.

\begin{figure}[h]
    \centering
    \includegraphics[width=0.85\textwidth]{III_1__scatter_batch_size.png}
    \caption{Pengaruh ukuran batch terhadap rata-rata \emph{latency} (Wu, 2020)}
    \label{fig:batch-latency}
\end{figure}

Untuk memahami perilaku sistem tersebut, beberapa variabel dapat dianggap sebagai representasi kondisi atau state sistem.
Backlog event menjadi indikator utama yang menggambarkan selisih antara laju produksi dan laju konsumsi pada saat itu.
Latensi pemrosesan event menunjukkan waktu yang dibutuhkan untuk memproses sebuah perubahan dan sangat dipengaruhi oleh besarnya batch, kapasitas komputasi, serta tingkat kontensi sumber daya.
Selain itu, utilisasi CPU atau memori pada komponen konsumer juga merefleksikan kapasitas sistem dalam merespons beban tambahan.
Variabel-variabel ini merepresentasikan kondisi sesaat sekaligus menentukan apakah sistem akan stabil atau terus menumpuk backlog.

Di sisi lain, variabel yang dapat dikendalikan atau \emph{control variables} meliputi ukuran batch dan interval polling pada konsumer.
Kedua parameter ini secara langsung mengatur ritme konsumsi data dan besarnya unit kerja yang diproses dalam setiap siklus.
Karena keduanya memengaruhi backlog, latensi, dan utilisasi sumber daya, kombinasi nilai yang tidak tepat dapat menyebabkan kondisi sistem berada di luar rentang yang diinginkan.
Masalah inilah yang ingin diatasi melalui pendekatan pengendalian dengan menentukan nilai pengaturan yang dapat menyeimbangkan pergerakan variabel state agar tetap berada pada kondisi stabil meskipun terjadi variasi beban.

Namun hubungan antara variabel state dan kontrol bersifat dinamis dan tidak selalu mudah dipetakan secara langsung.
Backlog, misalnya, bergantung pada interaksi antara laju produksi event dan kapasitas pemrosesan aktual pada konsumer, yang keduanya dapat bervariasi seiring waktu.
Latensi pemrosesan batch dapat berubah akibat peningkatan jumlah thread aktif, perubahan beban komputasi, atau efek kontensi I/O.
Ketidakpastian dan variasi inilah yang membuat pengaturan statis tidak memadai, karena parameter yang optimal pada satu kondisi dapat menjadi suboptimal atau bahkan kontraproduktif pada kondisi lainnya.
Oleh karena itu, diperlukan model yang mampu menangkap hubungan dinamis antara state dan kontrol sehingga sistem dapat diatur secara adaptif sesuai perubahan kondisi.

Analisis ini menunjukkan bahwa inti permasalahan bukan sekadar menentukan nilai batch atau interval polling yang paling cepat atau paling efisien, tetapi menemukan cara untuk menjaga keseimbangan antara backlog, latensi, dan utilisasi sumber daya dalam kondisi sistem yang terus berubah.
Model pengendalian dibutuhkan untuk menjaga agar variabel state bergerak menuju kondisi yang stabil dan terkendali, terlepas dari variasi beban.
Hasil analisis ini menjadi dasar perancangan pendekatan pengendalian pada bagian berikutnya.

\section{Analisis Solusi}

Berdasarkan analisis pada bagian sebelumnya, permasalahan sinkronisasi data pada arsitektur CQRS menunjukkan karakteristik sistem dinamis yang dipengaruhi oleh interaksi antara kondisi sistem dan parameter pengaturan pemrosesan.
Backlog, latensi sinkronisasi, dan utilisasi sumber daya tidak berkembang secara independen, melainkan saling mempengaruhi seiring perubahan beban dan konfigurasi sinkronisasi.
Oleh karena itu, solusi yang dirancang dalam penelitian ini tidak bertujuan mencari konfigurasi parameter statis yang optimal, melainkan mekanisme pengendalian yang mampu menyesuaikan parameter sinkronisasi secara adaptif agar sistem tetap berada pada kondisi seimbang.

Pendekatan pengendalian yang diusulkan dirancang dengan memformulasikan proses sinkronisasi sebagai sistem dinamis dalam kerangka \emph{state-space}.
Dengan formulasi ini, hubungan antara kondisi sistem (\emph{state}) dan parameter yang dapat dikendalikan (\emph{control input}) dapat dianalisis dan diatur secara sistematis.
Pendekatan ini memungkinkan penerapan berbagai strategi pengendalian, mulai dari kontrol umpan balik konvensional hingga kontrol optimal dan pendekatan berbasis data, dalam kerangka yang konsisten.

\subsection{Abstraksi Sistem dan Pemodelan State-Space}

Untuk merancang mekanisme pengendalian, proses sinkronisasi CQRS diabstraksikan sebagai sistem dinamis diskrit yang berevolusi seiring waktu.
Abstraksi ini memandang pipeline sinkronisasi, mulai dari konsumsi event hingga materialisasi \emph{read model}, sebagai sebuah sistem yang memiliki keadaan internal dan dapat dipengaruhi melalui aksi kontrol tertentu.
Pendekatan ini sejalan dengan praktik pemodelan sistem komputasi sebagai sistem dinamis ketika perilaku sistem dipengaruhi oleh umpan balik dan perubahan beban (Ogata, 2010).

Pada message sink, nilai-nilai matriks A, B, C, dan D diperoleh dari eksperimen dengan pendekatan \emph{open-loop} agar diketahui pengaruh variabel-variabel kontrol terhadap variabel-variabel state.
Eksperimen \emph{open-loop} adalah eksperimen yang menjalankan sistem tanpa kontroler aktif, sehingga parameter input divariasikan secara manual dan respons sistem dicatat apa adanya.
Lawannya adalah eksperimen \emph{closed-loop} yang menjalankan kontroler aktif sehingga parameter input ditentukan oleh kontroler berdasarkan kondisi sistem saat itu.
Vektor aksi kontrol didefinisikan sebagai:

\begin{equation}
\mathbf{u} = \begin{bmatrix} b \\ p \end{bmatrix}
\label{eq:vektor-kontrol}
\end{equation}

di mana:
\begin{enumerate}
\item $b$ adalah \emph{batch size}, yaitu jumlah pesan yang diproses dalam satu batch
\item $p$ adalah \emph{poll interval}, yaitu waktu jeda yang ditentukan untuk sistem dapat memproses batch data
\end{enumerate}

Vektor state $\mathbf{x}$ didefinisikan sebagai:

\begin{equation}
\mathbf{x} = \begin{bmatrix} n \\ c \\ m \\ o \end{bmatrix}
\label{eq:vektor-state}
\end{equation}

di mana:
\begin{enumerate}
\item $n$ adalah banyaknya pesan pada event source
\item $c$ adalah \emph{CPU utilization}
\item $m$ adalah \emph{memory utilization}
\item $o$ adalah \emph{I/O operation}
\end{enumerate}

Latensi sinkronisasi $l$ diperlakukan sebagai variabel observasi yang masuk ke matriks
output $C$ pada model state-space, bukan sebagai komponen vektor state.
Artinya, $l$ diukur untuk mengevaluasi kualitas sinkronisasi tetapi tidak digunakan oleh
kontroler untuk menentukan aksi koreksi.
Latensi diukur secara langsung melalui mekanisme \emph{manual offset commit} pada
\emph{message broker}.
Consumer tidak melakukan \emph{auto-commit} offset setelah mengambil batch, melainkan menunda
commit hingga \emph{read database} mengirimkan \emph{acknowledgement} yang mengonfirmasi
seluruh data dalam batch berhasil ditulis.
Latensi sinkronisasi untuk setiap batch dihitung sebagai:

\begin{equation}
l = t_{\text{ack}} - t_{\text{produce}}
\label{eq:latency-measurement}
\end{equation}

\noindent di mana $t_{\text{produce}}$ adalah waktu event diproduksi pada sisi tulis dan
$t_{\text{ack}}$ adalah waktu \emph{acknowledgement} diterima dari \emph{read database}.
Pendekatan ini memastikan latensi yang diukur mencakup seluruh waktu propagasi end-to-end,
termasuk waktu tunggu di antrian dan waktu pemrosesan batch.

Pemilihan variabel state didasarkan pada dua pertimbangan utama.
Pertama, variabel dapat diobservasi atau diukur secara real-time.
Kedua, variabel tersebut memiliki pengaruh langsung terhadap laju sinkronisasi dan tekanan
terhadap \emph{read database}.
Latensi tidak dimasukkan ke vektor state karena pengujian awal menunjukkan bahwa latensi di
state menyebabkan kontroler menekan \emph{batch size} secara berlebihan, sehingga latensi
hanya dijadikan variabel observasi terpisah.

Dalam konteks sistem nyata, hubungan antara \emph{state} dan \emph{control} bersifat nonlinier.
Namun, untuk keperluan perancangan kontrol optimal, sistem diasumsikan dapat didekati secara linier pada rentang operasi tertentu (\emph{linearization around operating point}).
Selain itu, diasumsikan bahwa perubahan beban dalam satu interval kontrol relatif kecil (\emph{quasi-stationary assumption}), sehingga pendekatan linier diskrit masih valid untuk perancangan kendali (Ogata, 2010; Hellerstein dkk., 2004).
Validitas asumsi ini kemudian diuji melalui eksperimen \emph{open-loop} pada tahap persiapan penelitian.

\subsection{Perancangan Fungsi Biaya dan Definisi Keseimbangan}

Tujuan pengendalian dalam penelitian ini adalah mengompromikan kecepatan sinkronisasi data dengan efisiensi penggunaan sumber daya.
Kompromi ini tidak dapat diwakili oleh satu metrik tunggal, karena melibatkan beberapa tujuan yang saling bertentangan, seperti minimisasi backlog, pengurangan latensi sinkronisasi, dan pembatasan utilisasi sumber daya agar tidak melampaui kapasitas sistem (Bailis dkk., 2012).

Untuk merepresentasikan tujuan tersebut secara formal, digunakan fungsi biaya kuadratik:

\begin{equation}
J = \sum_{t=0}^{\infty} \left( x_t^T Q x_t + u_t^T R u_t \right)
\label{eq:fungsi-biaya}
\end{equation}

Pada setiap langkah waktu $t$, vektor state $x_t$ berisi kondisi aktual sistem, yaitu jumlah backlog, utilisasi CPU, utilisasi memori, dan intensitas operasi I/O saat itu.
Masing-masing komponen $x_t$ dikuadratkan dan dikalikan dengan bobot yang ditetapkan di matriks $Q$.
Jika backlog saat ini adalah 3.000 dan target adalah 0, maka backlog berkontribusi pada biaya sebesar $3000^2 \times Q[\text{backlog}]$.
Hal yang sama berlaku untuk setiap variabel state lainnya.
Secara bersamaan, aksi kontrol yang diterapkan ($u_t$, yaitu batch size dan poll interval yang dipilih) juga dihitung biayanya melalui matriks $R$, agar sistem tidak mengambil tindakan yang terlalu ekstrem.
Nilai dari kedua komponen ini dijumlahkan untuk setiap langkah waktu, dan total akumulasinya itulah yang membentuk $J$.
Komponen pertama dari fungsi biaya dengan demikian memberikan penalti terhadap deviasi kondisi sistem dari keadaan yang diinginkan, sementara komponen kedua memberikan penalti terhadap besarnya aksi kontrol yang diterapkan.
Dengan formulasi ini, sistem diarahkan untuk mengurangi backlog dan utilisasi sumber daya sekaligus menghindari perubahan parameter sinkronisasi yang terlalu agresif maupun terlalu lambat.

Makna optimal dalam kerangka LQR didefinisikan sebagai keputusan kontrol yang meminimalkan fungsi biaya tersebut sepanjang horizon waktu, sehingga menghasilkan kompromi terbaik antara perbaikan kondisi sistem dan besarnya perubahan parameter sinkronisasi.
Bobot pada matriks $Q$ mencerminkan tingkat kepentingan relatif masing-masing variabel state, misalnya penekanan yang lebih besar pada backlog akan mendorong sistem untuk mempercepat sinkronisasi, sementara penekanan pada utilisasi sumber daya akan menghasilkan perilaku pengendalian yang lebih konservatif.
Matriks $R$ berfungsi untuk membatasi perubahan \emph{batch size} dan \emph{poll interval} agar tidak terlalu agresif, sehingga mengurangi potensi osilasi dan ketidakstabilan sistem.

\subsection{Kontrol Optimal Linear Quadratic Regulator (LQR)}

Kontrol optimal LQR digunakan sebagai pendekatan berbasis model yang secara eksplisit memformulasikan tujuan pengendalian dalam bentuk fungsi biaya.
LQR bekerja berdasarkan representasi \emph{state-space} sistem, di mana kondisi sinkronisasi direpresentasikan sebagai vektor state dan parameter sinkronisasi sebagai vektor kontrol.
Dengan model ini, LQR menghitung kebijakan kontrol yang meminimalkan fungsi biaya kuadratik yang merepresentasikan keseimbangan antara perbaikan kondisi sistem dan besarnya aksi kontrol.

Solusi optimal dari fungsi biaya ini adalah hukum kontrol linier:

\begin{equation}
u_t = -K x_t
\label{eq:hukum-kontrol}
\end{equation}

dengan $K$ adalah matriks penguat (gain) yang diperoleh dari penyelesaian persamaan Riccati.
Matriks ini menunjukkan seberapa besar tindakan korektif yang harus diambil untuk setiap kesalahan yang terdeteksi dalam status sistem.

LQR mempertimbangkan seluruh variabel state secara simultan dan tidak hanya berfokus pada satu sinyal kesalahan.
Pendekatan ini memungkinkan pengendalian yang lebih sistematis terhadap trade-off antara backlog, latensi sinkronisasi, dan utilisasi sumber daya.
Kebijakan kontrol LQR diperoleh melalui penyelesaian persamaan Riccati, sehingga aksi kontrol yang dihasilkan bersifat optimal secara matematis pada asumsi linearitas lokal dan gangguan terbatas.

Keterbatasan utama LQR terletak pada kebutuhan akan model sistem dan penentuan bobot fungsi biaya yang tepat.
Proses identifikasi model dan penalaan bobot memerlukan eksperimen awal dan pemahaman terhadap dinamika sistem.
Meskipun demikian, LQR memberikan kerangka pengendalian yang interpretable dan konsisten secara teori, sehingga cocok digunakan sebagai pendekatan utama dalam penelitian ini.

\subsection{Pendekatan Pengendalian yang Dibandingkan}

Untuk mengevaluasi efektivitas LQR secara menyeluruh, penelitian ini membandingkan empat pendekatan pengendalian yang merepresentasikan spektrum dari tanpa kontrol hingga kontrol optimal berbasis model.
Keempat pendekatan ini (statik, berbasis aturan, PID, dan LQR) menggunakan abstraksi sistem dan fungsi biaya yang sama sebagaimana didefinisikan pada Persamaan~\ref{eq:vektor-kontrol}--\ref{eq:vektor-state} dan Persamaan~\ref{eq:fungsi-biaya}, sehingga perbandingan dilakukan pada kerangka evaluasi yang konsisten.

\subsubsection{Kontrol Statik (Fixed)}

Pendekatan kontrol statik menggunakan nilai batch size dan poll interval yang tetap sepanjang eksperimen.
Nilai parameter ditentukan berdasarkan rata-rata kondisi operasional normal yang diperoleh dari eksperimen open-loop.
Pendekatan ini berfungsi sebagai baseline untuk menunjukkan kebutuhan akan mekanisme pengendalian adaptif.
Tanpa kemampuan menyesuaikan parameter terhadap perubahan beban, kontrol statik diharapkan menunjukkan kinerja yang menurun secara signifikan ketika kondisi operasional menyimpang dari asumsi awal.

\subsubsection{Kontrol Berbasis Aturan (Rule-Based)}

Pendekatan berbasis aturan menggunakan sekumpulan aturan threshold untuk menyesuaikan parameter sinkronisasi.
Aturan dirancang berdasarkan pemahaman domain tentang hubungan antara variabel state dan parameter kontrol.
Sebagai contoh, ketika backlog melebihi batas tertentu, batch size diperbesar untuk meningkatkan laju konsumsi; ketika utilisasi CPU melebihi ambang yang ditentukan, batch size dikurangi untuk menghindari kelebihan beban.

Pendekatan ini merepresentasikan praktik umum dalam pengelolaan sistem perangkat lunak, di mana administrator menetapkan aturan penyesuaian berdasarkan pengalaman operasional.
Keterbatasan utama pendekatan ini adalah ketidakmampuannya mempertimbangkan interaksi antar variabel secara simultan dan tidak adanya definisi formal mengenai keseimbangan yang ingin dicapai.

\subsubsection{Kontrol PID}

Kontrol PID diterapkan sebagai pendekatan kontrol klasik yang menggunakan satu sinyal kesalahan untuk menghasilkan aksi koreksi.
Dalam penelitian ini, PID dikonfigurasi sebagai kontroler single-loop dengan backlog sebagai variabel yang dikendalikan dan batch size sebagai variabel kontrol.
Sinyal kesalahan dihitung sebagai selisih antara nilai backlog aktual dan nilai referensi yang diinginkan.

Parameter $K_p$, $K_i$, dan $K_d$ ditentukan melalui metode tuning berdasarkan data eksperimen open-loop.
PID berfungsi sebagai representasi pendekatan kontrol klasik yang sudah mapan, dengan keterbatasan utama pada sifat SISO yang hanya mengamati satu variabel state dan mengendalikan satu variabel kontrol, tanpa memperhitungkan trade-off multi-variabel secara eksplisit.

\section{Rancangan Solusi}

Persiapan penelitian dilakukan untuk memperoleh data yang merepresentasikan dinamika sinkronisasi sistem pada berbagai kondisi beban.
Data ini digunakan untuk identifikasi hubungan antara kondisi sistem dan parameter sinkronisasi, serta sebagai dasar pelatihan model pembelajaran berbasis data.
Seluruh proses persiapan dirancang agar konsisten dengan tujuan pengendalian dan metrik evaluasi yang digunakan pada tahap pengujian.

\subsection{Eksperimen Open-Loop dan Identifikasi Sistem}

Eksperimen open-loop dilakukan dengan menjalankan sistem sinkronisasi tanpa mekanisme pengendalian adaptif aktif.
Parameter sinkronisasi diatur secara manual dan divariasikan dalam rentang yang telah ditentukan, sementara kondisi sistem diamati dan dicatat secara periodik.
Pendekatan ini dipilih untuk memastikan bahwa hubungan sebab akibat antara parameter sinkronisasi dan respons sistem dapat diamati secara langsung, tanpa dipengaruhi oleh umpan balik pengendali.

Pada setiap interval pengamatan, kondisi sistem dicatat sebagai vektor state yang mencakup backlog event, utilisasi CPU, utilisasi memori, dan intensitas operasi I/O pada read database.
Secara bersamaan, parameter kontrol berupa batch size dan poll interval dicatat sebagai bagian dari konfigurasi eksperimen.
Pengukuran dilakukan pada interval waktu yang tetap untuk menjaga konsistensi temporal antar variabel.

Eksperimen dijalankan pada beberapa skenario beban yang berbeda, termasuk beban rendah, sedang, dan tinggi.
Untuk setiap skenario, beberapa kombinasi parameter sinkronisasi diuji dan dijalankan selama jendela waktu observasi yang sama.
Selama periode tersebut, metrik performa seperti perubahan backlog, latensi sinkronisasi, dan utilisasi sumber daya dicatat untuk dianalisis lebih lanjut.
Data tersebut dimanfaatkan untuk mengamati karakteristik dinamika sistem dan merancang matriks kontrol.

\subsection{Definisi Aksi Kontrol Optimal}

Aksi kontrol optimal $u^*(t)$ didefinisikan sebagai parameter sinkronisasi yang menghasilkan kinerja terbaik untuk suatu kondisi sistem tertentu berdasarkan tujuan pengendalian yang telah ditetapkan.
Dalam penelitian ini, definisi optimal tidak diartikan sebagai minimisasi satu metrik secara terpisah, melainkan sebagai pencapaian keseimbangan antara backlog, latensi sinkronisasi, dan pemanfaatan sumber daya dalam suatu horizon pengamatan yang terbatas.

Secara formal, untuk kondisi sistem $x(t)$, aksi kontrol optimal didefinisikan sebagai:

\begin{equation}
u^*(t) = \arg\min_{u \in \mathcal{U}} J(u, x(t))
\label{eq:aksi-optimal}
\end{equation}

di mana $\mathcal{U}$ adalah himpunan kandidat aksi kontrol yang memenuhi batas operasional sistem, dan $J$ adalah fungsi biaya yang merepresentasikan tujuan pengendalian sesuai dengan Persamaan~\ref{eq:fungsi-biaya}.

Pada suatu kondisi awal sistem yang diamati, beberapa kombinasi parameter sinkronisasi, yaitu batch size dan poll interval, diuji secara terpisah.
Selama satu periode pengamatan, aksi kontrol dijaga tetap konstan untuk setiap kombinasi yang diuji, sementara respons sistem dicatat secara periodik.
Nilai fungsi biaya kemudian dihitung dari respons tersebut, dan kombinasi parameter yang menghasilkan nilai fungsi biaya minimum ditetapkan sebagai aksi kontrol optimal untuk kondisi sistem tersebut.

Pendekatan ini memastikan bahwa definisi $u^*(t)$ konsisten dengan tujuan pengendalian LQR, meskipun diperoleh melalui evaluasi empiris.
Dengan demikian, $u^*(t)$ menjadi aksi kontrol paling seimbang untuk kondisi sistem tertentu dalam batas asumsi dan horizon pengamatan yang berlaku, dan menjadi acuan pada analisis performa serta perbandingan di tahap evaluasi.

\section{Desain Eksperimen dan Skenario Pengujian}

Desain eksperimen pada penelitian ini bertujuan untuk mengevaluasi efektivitas pengendalian sinkronisasi data berbasis LQR dalam berbagai kondisi operasional sistem.
Eksperimen dirancang agar hasil yang diperoleh mencerminkan pengaruh mekanisme pengendalian, bukan faktor eksternal lainnya.

Eksperimen dilakukan dengan menjalankan sistem CQRS pada lingkungan uji yang konsisten, dengan konfigurasi perangkat keras, skema data, dan logika materialisasi read model yang dijaga tetap.
Variasi yang diperkenalkan dalam eksperimen dibatasi pada skenario beban dan parameter pengendalian yang menjadi fokus penelitian.
Pendekatan ini memastikan validitas internal eksperimen dan memungkinkan interpretasi hasil yang lebih jelas.

\subsection{Lingkungan Uji}

Lingkungan uji dirancang untuk merepresentasikan kondisi operasional sistem sinkronisasi data berbasis CQRS secara realistis dan terkontrol.
Seluruh eksperimen dijalankan pada arsitektur yang sama untuk memastikan bahwa perbedaan kinerja yang diamati berasal dari mekanisme pengendalian, bukan dari perubahan konfigurasi sistem atau lingkungan eksekusi.

Arsitektur sistem uji terdiri dari empat komponen utama, yaitu write model sebagai sumber perubahan data, message broker sebagai media distribusi event, message sink sebagai konsumer event sekaligus titik penerapan kontrol, dan read database sebagai target materialisasi data.
Alur sinkronisasi dimulai dari operasi tulis pada write model yang menghasilkan event perubahan data, kemudian event tersebut dipublikasikan ke message broker dan dikonsumsi oleh message sink untuk diproses secara bertahap sebelum diterapkan ke read database.

\begin{figure}[H]
    \centering
    \includegraphics[width=0.9\textwidth]{III_4_1__arsitektur_uji.png}
    \caption{Arsitektur Desain yang akan diuji}
    \label{fig:arsitektur-desain}
\end{figure}

Mekanisme kontrol LQR ditempatkan pada message sink karena komponen ini memiliki peran langsung dalam menentukan laju konsumsi event dan intensitas beban yang diberikan ke read database.
Parameter sinkronisasi yang dikendalikan meliputi batch size dan poll interval, yang secara langsung memengaruhi jumlah event yang diproses dalam satu siklus serta frekuensi pemrosesan.
Dengan menempatkan kontrol pada message sink, sistem memperoleh fleksibilitas untuk menyesuaikan laju sinkronisasi tanpa mengganggu logika bisnis pada write model maupun pola akses pada sisi baca.

Selama eksperimen, konfigurasi perangkat keras dan perangkat lunak dijaga tetap.
Kapasitas CPU, memori, dan penyimpanan tidak diubah, demikian pula versi sistem operasi, message broker, dan basis data yang digunakan.
Skema data pada write model dan read database serta logika materialisasi data juga tidak mengalami perubahan.
Pendekatan ini dilakukan untuk menjaga validitas internal eksperimen dan memastikan bahwa dinamika sistem yang diamati mencerminkan pengaruh pengendalian terhadap proses sinkronisasi.

Sebelum setiap eksperimen dimulai, sistem dijalankan hingga mencapai kondisi awal yang stabil dengan konfigurasi parameter sinkronisasi default.
Periode stabilisasi ini bertujuan untuk menghindari pengaruh transien awal sistem terhadap hasil pengukuran.
Setelah kondisi awal tercapai, mekanisme kontrol LQR diaktifkan dan eksperimen dijalankan sesuai dengan skenario beban yang telah ditentukan.

\subsection{Perancangan Skenario Beban}

Skenario beban dirancang untuk mengevaluasi kinerja pengendalian sinkronisasi data berbasis LQR pada berbagai kondisi operasional dan pola perubahan beban.
Selain variasi tingkat beban, penelitian ini juga mempertimbangkan karakteristik temporal perubahan beban karena respons sistem kendali sangat dipengaruhi oleh bagaimana beban berubah terhadap waktu.

Skenario beban dengan pola \emph{step} digunakan untuk merepresentasikan perubahan beban yang terjadi secara tiba-tiba, misalnya lonjakan laju produksi event akibat peningkatan trafik secara mendadak.
Pola ini digunakan untuk mengamati respons transien sistem, khususnya waktu respons awal, kenaikan backlog, dan kemampuan pengendalian dalam meredam lonjakan beban secara cepat.

Pola \emph{ramp} digunakan untuk merepresentasikan perubahan beban yang meningkat atau menurun secara bertahap.
Pola ini mencerminkan kondisi pertumbuhan trafik yang progresif dan memungkinkan evaluasi kemampuan pengendalian dalam menyesuaikan parameter sinkronisasi secara gradual tanpa menimbulkan osilasi yang berlebihan.

Pola \emph{impulse} digunakan untuk merepresentasikan lonjakan beban sesaat dengan durasi singkat (\emph{short burst}), yang kemudian diikuti oleh kembalinya beban ke kondisi awal.
Pola ini bertujuan untuk mengamati sensitivitas sistem terhadap gangguan sementara dan kemampuan pengendalian dalam memulihkan kondisi sistem setelah gangguan terjadi.

Selain itu, digunakan pola \emph{step periodik} untuk merepresentasikan perubahan beban yang berulang antara dua tingkat tertentu dalam interval waktu tetap.
Pola ini digunakan untuk mengevaluasi stabilitas jangka menengah sistem pengendalian serta potensi munculnya osilasi akibat perubahan beban yang berulang.
Pengamatan pada skenario ini difokuskan pada konsistensi respons sistem dan kemampuan pengendalian dalam mempertahankan performa yang stabil pada siklus beban yang berulang.

Terakhir, pola \emph{step low} digunakan untuk merepresentasikan kondisi beban rendah yang konstan (100 msg/s), berbeda dengan pola \emph{step} yang menggunakan rate jauh lebih tinggi.
Pola ini berguna untuk menguji apakah kontroler tetap stabil saat antrean sangat kecil dan tidak bereaksi berlebihan ketika sistem sebenarnya dalam kondisi tenang.

Dengan menggunakan kombinasi kelima pola beban tersebut, eksperimen dapat menguji respons dan stabilitas pengendalian pada berbagai kondisi operasional.

\begin{figure}[H]
    \centering
    \begin{tabular}{cc}
        \includegraphics[width=0.48\textwidth]{III_4_2_a__skenario_step.png} &
        \includegraphics[width=0.48\textwidth]{III_4_2_b__skenario_ramp.png} \\
        (a) Fungsi Step & (b) Fungsi Ramp \\[0.5cm]
        \includegraphics[width=0.48\textwidth]{III_4_2_c__skenario_impulse.png} &
        \includegraphics[width=0.48\textwidth]{III_4_2_d__skenario_periodic.png} \\
        (c) Fungsi Impulse & (d) Fungsi Step Periodik
    \end{tabular}
    \caption[Simulasi beban dengan beberapa skenario fungsi]{Simulasi beban dengan beberapa skenario fungsi (a) fungsi step, (b) fungsi ramp, (c) fungsi impulse dan (d) fungsi step yang berulang}
    \label{fig:skenario-beban}
\end{figure}

\subsection{Perancangan Eksperimen Perbandingan Metode}

Untuk memastikan keadilan perbandingan, seluruh pendekatan pengendalian dijalankan pada skenario beban yang sama dengan kondisi awal sistem yang identik.
Setiap metode menerima informasi state yang sama pada setiap interval kontrol dan menghasilkan aksi kontrol sesuai mekanismenya masing-masing.
Interval pengambilan keputusan kontrol dijaga tetap untuk seluruh metode agar perbedaan kinerja yang diamati berasal dari strategi pengendalian, bukan dari perbedaan frekuensi pengambilan keputusan.

Konfigurasi masing-masing metode ditentukan sebagai berikut.
Kontrol statik menggunakan nilai batch size dan poll interval yang ditetapkan berdasarkan rata-rata parameter pada kondisi beban normal dari data eksperimen open-loop.
Kontrol berbasis aturan menggunakan sekumpulan aturan threshold yang dirancang berdasarkan hubungan antar variabel yang teramati pada eksperimen open-loop, dengan batas threshold ditentukan dari distribusi nilai variabel state selama eksperimen.
Kontrol PID dikonfigurasi sebagai kontroler single-loop pada variabel backlog, dengan parameter $K_p$, $K_i$, dan $K_d$ ditentukan melalui metode tuning berbasis data eksperimen open-loop.
Kontrol LQR menggunakan matriks $A$ dan $B$ yang diidentifikasi dari eksperimen open-loop, serta matriks $Q$ dan $R$ yang mencerminkan prioritas pengendalian yang telah ditetapkan pada Persamaan~\ref{eq:fungsi-biaya}.

Setiap metode dijalankan pada kelima skenario beban (step, ramp, impulse, step periodik, dan step low) dengan durasi eksperimen yang sama.
Hasil eksperimen dievaluasi menggunakan seluruh metrik yang telah didefinisikan pada bagian III.5, sehingga perbandingan mencakup aspek kinerja sinkronisasi, dinamika kontrol, efisiensi pemrosesan, dan biaya kontrol.

\section{Metrik Evaluasi Kerja}

Evaluasi kinerja pada penelitian ini menilai sejauh mana pengendali menyeimbangkan kualitas sinkronisasi data dengan efisiensi penggunaan sumber daya.
Oleh karena itu, metrik evaluasi dipilih secara selektif agar dapat merepresentasikan tujuan pengendalian yang telah dirumuskan pada tahap perancangan solusi, serta mampu menangkap respons sistem pada berbagai kondisi beban dan pola perubahan beban.

Metrik evaluasi dikelompokkan berdasarkan perannya dalam analisis kinerja sistem, meliputi metrik inti kinerja sinkronisasi, metrik dinamika kontrol, metrik efisiensi pemrosesan, dan metrik biaya kontrol.
Pengelompokan ini bertujuan untuk memisahkan pengukuran kualitas sinkronisasi, karakteristik respons dinamis, efisiensi operasional, dan kesesuaian solusi terhadap formulasi optimasi yang digunakan.
Seluruh metrik diukur secara periodik selama eksperimen dan digunakan secara konsisten sebagai dasar analisis dan perbandingan keempat pendekatan pengendalian pada bab hasil dan pembahasan.

\subsection{Metrik Inti Kinerja Sinkronisasi}

Kelompok metrik inti ini menilai kualitas dasar proses sinkronisasi data pada arsitektur CQRS.
Metrik ini dipilih karena secara langsung merepresentasikan tujuan utama pengendalian, yaitu menjaga konsistensi data dan kestabilan sistem pada berbagai kondisi beban.
Seluruh metrik inti diukur secara periodik selama eksperimen dan dianalisis baik pada kondisi tunak maupun transien.

Backlog sinkronisasi digunakan sebagai indikator utama untuk mengukur tingkat ketertinggalan proses sinkronisasi.
Backlog didefinisikan sebagai jumlah event yang telah dihasilkan oleh write model tetapi belum diproses oleh message sink pada suatu waktu tertentu.
Nilai backlog yang tinggi menunjukkan ketidakseimbangan antara laju produksi dan konsumsi event, sedangkan pengendalian yang efektif ditunjukkan oleh kemampuan menekan pertumbuhan backlog dan mempercepat pemulihannya setelah terjadi gangguan.
Backlog diukur dalam satuan jumlah event, dengan analisis mencakup nilai rata-rata, nilai maksimum, dan waktu pemulihan menuju kondisi referensi.

Latensi sinkronisasi digunakan untuk mengukur keterlambatan propagasi perubahan data dari write model hingga perubahan tersebut tersedia pada read database.
Latensi ini mencerminkan konsistensi data yang dapat diamati oleh pengguna dan berhubungan langsung dengan kualitas layanan sistem.
Pengukuran latensi dilakukan dengan mencatat selisih waktu antara timestamp event pada sisi tulis dan waktu penerapan perubahan pada sisi baca.
Satuan yang digunakan adalah milidetik atau detik, dengan analisis meliputi nilai rata-rata dan nilai maksimum selama periode eksperimen.

Utilisasi sumber daya digunakan untuk mengevaluasi dampak pengendalian terhadap kapasitas sistem, khususnya pada komponen read database yang menjadi target materialisasi data.
Metrik ini mencakup utilisasi CPU, penggunaan memori, dan aktivitas I/O selama proses sinkronisasi berlangsung.
Utilisasi CPU dan memori diukur dalam persen, sedangkan aktivitas I/O diukur dalam operasi per detik atau jumlah data yang diproses per satuan waktu.
Metrik ini digunakan untuk memastikan bahwa pengendalian mampu menjaga keseimbangan antara peningkatan laju sinkronisasi dan keterbatasan sumber daya sistem.

\begin{table}[H]
    \centering
    \caption{Metrik kinerja sinkronisasi}
    \label{tab:metrik-sinkronisasi}
    \renewcommand{\arraystretch}{1.4}
    \begin{tabular}{|c|p{3cm}|p{6cm}|p{2.5cm}|}
    \hline
    \textbf{No} & \textbf{Metrik} & \textbf{Definisi Operasional} & \textbf{Satuan} \\
    \hline
    1 & Backlog Sinkronisasi & Jumlah event yang telah dihasilkan oleh write model tetapi belum diproses oleh message sink pada waktu tertentu & Event \\
    \hline
    2 & Backlog Rata-rata & Nilai rata-rata backlog selama satu periode eksperimen & Event \\
    \hline
    3 & Backlog Maksimum & Nilai backlog tertinggi yang tercatat selama eksperimen & Event \\
    \hline
    4 & Waktu Pemulihan Backlog & Waktu yang dibutuhkan backlog untuk kembali mendekati nilai referensi setelah terjadi gangguan beban & Detik \\
    \hline
    5 & Latensi Sinkronisasi Rata-rata & Rata-rata waktu dari event dihasilkan hingga diterapkan pada read database & Milidetik/detik \\
    \hline
    6 & Latensi Sinkronisasi Maksimum & Nilai latensi sinkronisasi tertinggi selama eksperimen & Milidetik/detik \\
    \hline
    7 & Utilisasi CPU & Persentase penggunaan CPU pada read database selama proses sinkronisasi & Persen \\
    \hline
    8 & Utilisasi Memori & Persentase penggunaan memori pada read database selama proses sinkronisasi & Persen \\
    \hline
    9 & Aktivitas I/O & Intensitas operasi baca tulis selama proses sinkronisasi & Persen \\
    \hline
    \end{tabular}
\end{table}

\subsection{Metrik Dinamika Kontrol}

Kelompok metrik dinamika kontrol menilai karakteristik respons sistem terhadap perubahan beban, terutama pada kondisi transien seperti pola step, ramp, impulse, dan step periodik.
Berbeda dengan metrik inti yang menilai kualitas sinkronisasi secara umum, metrik dinamika kontrol berfokus pada bagaimana sistem bereaksi terhadap gangguan dan seberapa cepat sistem kembali ke kondisi yang diinginkan.
Metrik ini sangat relevan dalam konteks pengendalian berbasis LQR karena secara langsung mencerminkan kualitas respons kendali yang dihasilkan.

Evaluasi dinamika kontrol dilakukan dengan mengamati perilaku temporal backlog sebagai variabel state utama.
Respons backlog terhadap perubahan beban digunakan untuk menilai kecepatan respons, tingkat agresivitas pengendalian, dan kestabilan menuju kondisi referensi.
Dengan menggunakan metrik dinamika kontrol, penelitian ini dapat membedakan pengendalian yang cepat namun berisiko agresif dari pengendalian yang lebih stabil dan terkontrol.

Rise time digunakan untuk mengukur kecepatan respons pengendalian terhadap gangguan beban.
Nilai rise time yang kecil menunjukkan bahwa mekanisme kontrol mampu merespons perubahan kondisi sistem secara cepat.
Namun, respons yang terlalu cepat dapat meningkatkan risiko overshoot, sehingga rise time perlu dianalisis bersamaan dengan metrik lainnya.

Overshoot merepresentasikan tingkat agresivitas pengendalian setelah terjadi perubahan beban.
Overshoot yang besar menunjukkan bahwa pengendalian memberikan respons berlebihan, yang berpotensi meningkatkan tekanan terhadap sistem.
Dalam konteks sinkronisasi data, overshoot backlog yang tinggi dapat mengindikasikan risiko overload sementara pada read database.

Settling time digunakan untuk mengevaluasi kestabilan respons pengendalian.
Metrik ini menunjukkan berapa lama sistem membutuhkan waktu untuk kembali dan bertahan pada kondisi yang mendekati nilai referensi setelah gangguan terjadi.
Settling time yang singkat tanpa overshoot berlebihan menunjukkan karakteristik pengendalian yang stabil dan efektif.

\begin{table}[H]
    \centering
    \caption{Metrik dinamika kontrol}
    \label{tab:metrik-dinamika}
    \renewcommand{\arraystretch}{1.4}
    \begin{tabular}{|c|p{3cm}|p{6cm}|p{2.5cm}|}
    \hline
    \textbf{No} & \textbf{Metrik} & \textbf{Definisi Operasional} & \textbf{Satuan} \\
    \hline
    1 & Rise Time & Waktu yang dibutuhkan backlog untuk kembali mendekati nilai referensi setelah terjadi perubahan beban & Detik \\
    \hline
    2 & Overshoot & Selisih maksimum backlog yang melebihi nilai referensi setelah terjadi perubahan beban & Persen \\
    \hline
    3 & Settling Time & Waktu yang dibutuhkan hingga backlog berada dan bertahan dalam rentang toleransi tertentu di sekitar nilai referensi & Detik \\
    \hline
    \end{tabular}
\end{table}

\subsection{Metrik Efisiensi Pemrosesan}

Kelompok metrik efisiensi pemrosesan menilai sejauh mana pengendalian sinkronisasi meningkatkan kinerja pemrosesan event tanpa mengorbankan efisiensi penggunaan sumber daya.
Pada sistem CQRS, pengendalian yang efektif dicirikan oleh penurunan backlog dan latensi sekaligus kemampuan memproses event secara efisien dengan sumber daya yang ada.

Evaluasi efisiensi pemrosesan penting karena pengendalian yang terlalu agresif dapat meningkatkan throughput dalam jangka pendek, tetapi berpotensi menimbulkan pemborosan sumber daya atau tekanan berlebihan pada sistem.
Oleh karena itu, metrik efisiensi digunakan untuk menilai keseimbangan antara kapasitas pemrosesan yang dicapai dan biaya sumber daya yang dikeluarkan selama proses sinkronisasi berlangsung.

Throughput pemrosesan digunakan untuk mengukur kapasitas sistem dalam memproses event selama proses sinkronisasi.
Nilai throughput yang tinggi menunjukkan bahwa sistem mampu menangani laju event yang besar, namun metrik ini perlu dianalisis bersamaan dengan utilisasi sumber daya agar tidak menimbulkan interpretasi yang menyesatkan.

Efisiensi CPU dan efisiensi memori digunakan untuk mengevaluasi seberapa efektif sumber daya sistem dimanfaatkan dalam mendukung proses sinkronisasi.
Nilai efisiensi yang tinggi menunjukkan bahwa peningkatan throughput dicapai dengan penggunaan sumber daya yang relatif efisien.
Kombinasi metrik throughput dan efisiensi sumber daya membuat evaluasi kinerja mencakup kecepatan pemrosesan sekaligus keberlanjutan operasional sistem.

\begin{table}[H]
    \centering
    \caption{Metrik efisiensi pemrosesan}
    \label{tab:metrik-efisiensi}
    \renewcommand{\arraystretch}{1.4}
    \begin{tabular}{|c|p{3cm}|p{6cm}|p{2.5cm}|}
    \hline
    \textbf{No} & \textbf{Metrik} & \textbf{Definisi Operasional} & \textbf{Satuan} \\
    \hline
    1 & Throughput Pemrosesan & Jumlah event yang berhasil diproses oleh message sink per satuan waktu & Event/detik \\
    \hline
    2 & Efisiensi CPU & Rasio antara throughput pemrosesan dan utilisasi CPU & Event/detik/\% CPU \\
    \hline
    3 & Efisiensi Memori & Rasio antara throughput pemrosesan dan utilisasi memori & Event/detik/\% memori \\
    \hline
    \end{tabular}
\end{table}

\subsection{Metrik Biaya Kontrol}

Kelompok metrik biaya kontrol menilai kinerja pengendalian dari sudut pandang optimasi yang mendasari perancangan Linear Quadratic Regulator (LQR).
Berbeda dengan metrik kinerja dan efisiensi yang menilai dampak pengendalian pada sistem secara eksternal, metrik biaya kontrol menilai sejauh mana kebijakan kontrol yang diterapkan selaras dengan tujuan optimasi yang telah diformulasikan melalui fungsi biaya.

Evaluasi biaya kontrol penting untuk memastikan bahwa perbaikan kinerja sinkronisasi tidak dicapai melalui aksi kontrol yang berlebihan.
Dalam konteks ini, biaya kontrol digunakan untuk mengukur keseimbangan antara perbaikan kondisi sistem dan besarnya aksi kontrol yang diterapkan.
Dengan menganalisis nilai biaya kontrol, penelitian ini dapat menilai konsistensi perilaku pengendalian LQR terhadap tujuan optimasi yang dirancang.

Nilai fungsi biaya digunakan sebagai indikator utama untuk mengevaluasi performa pengendalian LQR dari perspektif optimasi.
Nilai fungsi biaya yang lebih rendah menunjukkan bahwa sistem berada lebih dekat dengan kondisi yang diinginkan dengan besaran aksi kontrol yang relatif terkendali.

Biaya state digunakan untuk menilai sejauh mana kondisi sistem menyimpang dari nilai referensi selama proses sinkronisasi.
Nilai biaya state yang tinggi mengindikasikan bahwa sistem sering berada jauh dari kondisi target, misalnya akibat backlog yang besar atau utilisasi sumber daya yang berlebihan.

Biaya aksi kontrol digunakan untuk mengevaluasi tingkat agresivitas pengendalian.
Nilai biaya aksi kontrol yang tinggi menunjukkan bahwa perbaikan kondisi sistem dicapai melalui perubahan parameter sinkronisasi yang besar atau sering.
Dengan menganalisis kedua komponen biaya secara terpisah, penelitian ini dapat mengevaluasi keseimbangan antara kualitas sinkronisasi dan intensitas aksi kontrol yang diterapkan.

\begin{table}[H]
    \centering
    \caption{Metrik biaya kontrol}
    \label{tab:metrik-biaya}
    \renewcommand{\arraystretch}{1.4}
    \begin{tabular}{|c|p{3cm}|p{6cm}|p{2.5cm}|}
    \hline
    \textbf{No} & \textbf{Metrik} & \textbf{Definisi Operasional} & \textbf{Satuan} \\
    \hline
    1 & Nilai Fungsi Biaya & Akumulasi penalti state dan aksi kontrol selama horizon pengamatan sesuai fungsi biaya kuadratik LQR & Tanpa dimensi \\
    \hline
    2 & Biaya State & Komponen fungsi biaya yang merepresentasikan penalti terhadap deviasi state dari nilai referensi & Tanpa dimensi \\
    \hline
    3 & Biaya Aksi Kontrol & Komponen fungsi biaya yang merepresentasikan penalti terhadap besarnya aksi kontrol yang diterapkan & Tanpa dimensi \\
    \hline
    \end{tabular}
\end{table}

\subsection{Analisis Sensitivitas Fungsi Biaya}

Nilai fungsi biaya $J$ bergantung pada pemilihan matriks pembobot $Q$ dan $R$ yang ditentukan oleh peneliti.
Pemilihan bobot yang berbeda akan menghasilkan nilai $J$ yang berbeda pula, sehingga kesimpulan yang diambil berdasarkan satu konfigurasi $Q$ dan $R$ tertentu berpotensi tidak berlaku secara umum.
Untuk mengatasi keterbatasan ini, dilakukan analisis sensitivitas dengan mengevaluasi nilai fungsi biaya seluruh metode pengendalian pada beberapa variasi konfigurasi $Q$ dan $R$.

Analisis sensitivitas dilakukan dengan mendefinisikan tiga konfigurasi pembobot yang merepresentasikan prioritas pengendalian yang berbeda.
Konfigurasi pertama memberikan bobot lebih besar pada komponen backlog sinkronisasi, sehingga merepresentasikan skenario di mana pengurangan backlog menjadi prioritas utama.
Konfigurasi kedua memberikan bobot lebih besar pada komponen utilisasi sumber daya (CPU dan I/O), sehingga merepresentasikan skenario di mana efisiensi penggunaan sumber daya menjadi perhatian utama.
Konfigurasi ketiga menggunakan bobot yang seimbang antara backlog dan sumber daya, sehingga merepresentasikan kompromi antara kedua tujuan.

Secara struktural, ketiga konfigurasi matriks pembobot $Q$ memiliki bentuk diagonal $4\times 4$ yang berkorespondensi dengan vektor state $\mathbf{x} = [n, c, m, o]^T$.
Pada konfigurasi prioritas backlog, elemen diagonal pertama (queue) dibobot jauh lebih tinggi dibanding elemen lainnya.
Pada konfigurasi prioritas sumber daya, elemen CPU dan I/O dibobot lebih tinggi sementara queue rendah.
Pada konfigurasi seimbang, queue dan sumber daya dibobot relatif setara.
Nilai bobot spesifik diturunkan menggunakan \emph{Bryson's rule} pada ruang ternormalisasi, dengan detail derivasi dan nilai numerik final ditampilkan pada Bab~IV.

Matriks $R$ ditetapkan konstan di seluruh konfigurasi sehingga variasi hanya terjadi pada bobot state, dan perbedaan nilai $J$ antar konfigurasi murni mencerminkan pengaruh prioritas pengendalian.

Pada setiap konfigurasi, nilai fungsi biaya $J$ dihitung untuk seluruh empat metode pengendalian, yaitu kontrol statik, rule-based, PID, dan LQR, menggunakan data trajectory yang sama dari hasil eksperimen. Dengan demikian, perbedaan nilai $J$ antar metode pada suatu konfigurasi mencerminkan perbedaan kinerja, bukan perbedaan data.

Untuk memperoleh ukuran perbandingan yang tidak bergantung pada skala absolut nilai $J$, digunakan \emph{normalized regret}. Pada setiap konfigurasi $k$, nilai fungsi biaya metode $i$ dinormalisasi terhadap metode dengan nilai $J$ terendah pada konfigurasi tersebut:

\begin{equation}
    \hat{J}_i^{(k)} = \frac{J_i^{(k)} - J_{\min}^{(k)}}{J_{\min}^{(k)}}
    \label{eq:normalized-regret}
\end{equation}

\noindent di mana $J_{\min}^{(k)} = \min_{i} J_i^{(k)}$ adalah nilai fungsi biaya terendah di antara seluruh metode pada konfigurasi $k$. Nilai $\hat{J}_i^{(k)} = 0$ menunjukkan bahwa metode $i$ adalah yang terbaik pada konfigurasi tersebut, sedangkan nilai yang lebih besar menunjukkan seberapa jauh metode tersebut dari performa terbaik secara proporsional.

Berdasarkan \emph{normalized regret}, adaptabilitas setiap metode terhadap variasi prioritas pengendalian dievaluasi menggunakan dua ukuran. Ukuran pertama adalah \emph{mean regret}, yang didefinisikan sebagai rata-rata \emph{normalized regret} di seluruh konfigurasi:

\begin{equation}
    \bar{\hat{J}}_i = \frac{1}{K} \sum_{k=1}^{K} \hat{J}_i^{(k)}
    \label{eq:mean-regret}
\end{equation}

\noindent di mana $K$ adalah jumlah konfigurasi pembobot yang diuji. Metode dengan \emph{mean regret} terendah dianggap sebagai metode terbaik secara keseluruhan karena secara rata-rata paling dekat dengan performa optimal di seluruh konfigurasi.

Ukuran kedua adalah \emph{max regret}, yang didefinisikan sebagai nilai \emph{normalized regret} tertinggi di antara seluruh konfigurasi:

\begin{equation}
    \hat{J}_i^{\max} = \max_{k} \hat{J}_i^{(k)}
    \label{eq:max-regret}
\end{equation}

\noindent Metode dengan \emph{max regret} terendah dianggap sebagai metode paling robust karena tidak pernah mengalami penurunan performa yang signifikan pada konfigurasi manapun.

Nilai \emph{mean regret} dan \emph{max regret} dihitung untuk seluruh empat metode pengendalian berdasarkan hasil evaluasi pada ketiga konfigurasi pembobot. Metode dengan \emph{mean regret} terendah dianggap sebagai metode terbaik secara keseluruhan karena secara rata-rata paling dekat dengan performa optimal, sedangkan metode dengan \emph{max regret} terendah dianggap sebagai metode paling robust karena tidak pernah mengalami degradasi performa yang signifikan pada konfigurasi manapun. Kesimpulan perbandingan metode dianggap robust apabila ranking relatif antar metode tetap konsisten di seluruh variasi konfigurasi $Q$ dan $R$. Apabila kedua ukuran menunjuk pada metode yang sama, maka kesimpulan keunggulan metode tersebut menjadi semakin kuat. Sebaliknya, apabila ranking berubah secara signifikan pada konfigurasi tertentu, maka keunggulan suatu metode dinyatakan bergantung pada prioritas pengendalian yang dipilih, dan hal ini dilaporkan sebagai bagian dari analisis hasil.
